\section{Detailed proofs and guide to the results}

These appendices give definitions, exact statements and analytic proofs for
all results summarized in the main text.  They also provide a reader's guide
to the specializations for finite contractions, finite and countable
pure-point spectra, observer families, local spin systems and finite-record
inference.  All arguments are deterministic.  No random-matrix premise,
numerical experiment or unproved nonresonance assumption is used.

\section{Carrier selection and arbitrary spectral type}

\subsection{Finite contractions}

Let $T$ be a contraction on a finite-dimensional complex Hilbert space.
Every eigenvalue satisfies $|\lambda|\le1$.  A Jordan block of size greater
than one on $|\lambda|=1$ would make $\|T^n\|$ grow polynomially, contrary
to $\|T^n\|\le1$.  Hence the peripheral part is semisimple.  Let
$P_\lambda$ be its spectral projections and put
\[
 x_{\rm per}=\sum_{|\lambda|=1}P_\lambda x,
 \qquad M_{\rm per}(T,x)=\sum_{|\lambda|=1}\|P_\lambda x\|^2.
\]

\begin{theorem}[Exact returnability transition]
For every finite contraction and vector,
\[
 \langle x,T^nx\rangle
 =\sum_{|\lambda|=1}\lambda^n\|P_\lambda x\|^2+r_n,
 \qquad r_n\to0.
\]
For every $a<M_{\rm per}(T,x)$ and every burn-in $B$, some $n>B$
satisfies $a\le|\langle x,T^nx\rangle|$.  For every
$a>M_{\rm per}(T,x)$, all sufficiently large $n$ satisfy
$|\langle x,T^nx\rangle|<a$.
\end{theorem}

\begin{proof}
The generalized eigenspaces with $|\lambda|<1$ decay, giving $r_n\to0$.
The peripheral sum is a positive weighted finite phase packet.  Simultaneous
Diophantine approximation returns all its phases arbitrarily close to one
after every delay, so its modulus approaches the sum of its weights,
$M_{\rm per}$, arbitrarily late.  The triangle inequality gives the upper
ceiling; the vanishing transient turns it into an eventual strict bound above
the ceiling.
\end{proof}

If at least two peripheral modes have positive mass, normalize their weights
by $M_{\rm per}$.  Applying the water deadline of Appendix~\ref{app:water} and waiting
until $|r_n|$ is below half the target slack gives an explicit post-burn-in
return below every strict ceiling target.  The same decomposition and
threshold hold termwise for every finite convex mixture of contractions.

\subsection{Stationary spectral measures and collision hierarchy}

For a probability measure $\mu$ on $\R$, let
$\widehat\mu(t)=\int\ee^{\ii tE}\,d\mu(E)$ and let
$\check\mu(A)=\mu(-A)$.  Its difference measure is
$\nu=\mu*\check\mu$, so
$|\widehat\mu(t)|^2=\widehat\nu(t)$.
Write $G_\mu$ for the additive subgroup of $\R$ generated by
$\{E-E':E,E'\in\operatorname{supp}\mu\}$; this is the visible
frequency-difference group.

\begin{theorem}[Stationary spectral return law]
For every $r\ge1$,
\[
 \lim_{T\to\infty}\frac1T\int_0^T
   |\widehat\mu(t)|^{2r}\,dt
 =\nu^{*r}(\{0\}).
\]
Every continuous statistic of $|\widehat\mu(t)|^2\in[0,1]$ has a
long-time mean.  At $r=1$, the mean is
$\sum_E\mu(\{E\})^2$.  Thus the stationary recurrent fidelity is exactly
the point-spectrum collision mass.
\end{theorem}

\begin{proof}
Expand the (2r)-fold product as the Fourier transform of $\nu^{*r}$.
Ces\`aro averaging sends $\ee^{\ii t\omega}$ to the indicator of
$\omega=0$; bounded convergence gives the formula.  Polynomial
approximation on $[0,1]$ extends convergence from moments to every
continuous statistic.
\end{proof}

\begin{theorem}[Pure-point, cyclic and atomless trichotomy]
\leavevmode
\begin{enumerate}[leftmargin=1.6em]
\item Every nonnegative summable countable phase energy
  $\sum_j a_j|\ee^{\ii t\omega_j}-1|^2$ falls below every positive
  resolution at some positive integer time, quantitatively by synchronizing
  a finite head and bounding the summable tail.
\item A probability spectrum has an exact unit-fidelity return at some
  positive time if and only if its visible frequency-difference group is
  cyclic.
\item If $\mu$ is atomless, its projective observer temporal dimension is
  continuum and every fixed positive-fidelity return window has asymptotic
  density zero.
\end{enumerate}
\end{theorem}

The first clause is the countable extension of simultaneous phase recurrence.
The second follows because equality in the triangle inequality forces all
visible phases to coincide.  The third combines atomless difference spectra
with the vanishing Ces\`aro collision mass; the continuum dimension is the
cardinal rank of the atomless phase representation.

\section{Finite projective completion and exact conservation}

Let $I$ be a finite $D$-element carrier and
\[
 P_I=\{k\in\Z^I:\sum_ik_i=0\},\qquad
 D_E(k)=\sum_i k_iE_i,
 \qquad\Rel_E=\ker D_E.
\]
Choose $i_0\in I$ and write
$G_E=\operatorname{span}_{\Q}\{E_i-E_{i_0}\}$.

\begin{theorem}[Projective recurrence completion]
The closure of the relative phase orbit is
\[
 \Hull_E=\Rel_E^\perp
 =\{z\in\T^I/\T:z^k=1\text{ for every }k\in\Rel_E\}.
\]
Haar measure on this hull is the completed phase law.  A balanced character
is almost surely one exactly when it belongs to $\Rel_E$.
\end{theorem}

\begin{proof}
Every relation annihilates the orbit and hence its closure.  Conversely, a
character that annihilates the orbit has real detuning zero, so it lies in
$\Rel_E$.  Equality follows from annihilator duality for compact abelian
groups.
\end{proof}

Tensoring the detuning map with $\Q$ gives an exact sequence
\[
 0\to\ker D_E\to P_I\otimes\Q\xrightarrow{D_E}G_E\to0.
\]
Since $\dim_\Q(P_I\otimes\Q)=D-1$, rank--nullity proves
\[
 \tau(E)+\rho(E)=D-1.
\]
The dimension of $\Hull_E$ is $\tau(E)$.

\subsection{Replica height and phase-grid readout}

Write $k=k^+-k^-$ with disjoint nonnegative parts.  Balancedness gives
$|k^+|_1=|k^-|_1=:h(k)$.  Expanding the $r$th fidelity replica shows that
the character $k$ has positive coefficient exactly when $h(k)\le r$.
Therefore its first possible replica order is $h(k)$, and a replica detects
it exactly when $k\in\Rel_E$ and $h(k)\le r$.

Each order-$r$ phase polynomial has coordinate Fourier bandwidth at most
$r$.  The discrete Fourier transform on a $(2r+1)$-point grid in each
cycled coordinate therefore recovers every coefficient exactly.  A smaller
one-coordinate grid aliases two admissible frequencies, proving worst-case
sharpness.  This proves the replica-height law, exact finite-grid
spectroscopy, and its Nyquist boundary.

\section{Coupling walls and conversion}

For countable profiles $E,W$, define the affine detuning
\[
 D_{E+gW}(k)=D_E(k)+gD_W(k).
\]
If $D_W(k)=0$, the character is either protected for every $g$ or never
crosses.  If $D_W(k)\ne0$, it crosses at the single strength
$-D_E(k)/D_W(k)$.  The character group is countable.

\begin{theorem}[Protected-relation intersection law]
There is a countable $C\subset\R$ such that for every $g\notin C$,
\[
 \Rel_{E+gW}=\Rel_E\cap\Rel_W.
\]
Every open interval contains a nonexceptional strength.  At a transverse
wall, the finite-window character reporter equals
$\sinc(TD_{E+gW}(k))$, with the exact inverse-time crossover profile.
\end{theorem}

The finite completed indicator is discontinuous at the wall, whereas the
reporter is continuous at every finite $T$.  Since arbitrarily small
frequency perturbations can have any prescribed rational rank, every
positive finite temporal-rank stratum is dense and the rank map is nowhere
continuous.  Applying Duhamel's identity to normalized observables gives
all-rank finite-horizon indistinguishability.  Intersecting over an allowed
perturbation space proves that robust relations are precisely
$\Rel_E$ intersected with its annihilator.

\subsection{Complete finite conversion}

Let spectra have $m+1$ and $n+1$ levels.  Their uncoupled product has at
least $mn$ rectangle relations.  A separable diagonal interaction can be
chosen, inside every positive operator-norm budget, so that the coupled
$(m+1)(n+1)$-level spectrum is maximally projective.  Consequently
\begin{align*}
 \rho_{\rm coupled}&=0,\\
 \tau_{\rm coupled}-\tau_{\rm uncoupled}&=\rho_{\rm uncoupled}\ge mn.
\end{align*}
Every positive full-support replica collision moment then equals its
unavoidable permutation floor.  The same interaction obeys the all-probe
bound
\[
 \|U_g(t)OU_g(t)^*-U_0(t)OU_0(t)^*\|
 \le2|t|\,\|H_g-H_0\|\,\|O\|.
\]
The maximal set in the finite separable parameter space is dense, residual
and full measure; its complement remains null in every positive-volume
parameter region.  Combining rational independence with torus
equidistribution gives the deterministic strict/full Haar-frequency
sandwich and the exact uniform-target half-dimensional volume exponent.

\subsection{Countable relation hull and weak recurrence ray}

For a countable pure-point carrier, finite characters still form a countable
group.  Its relation lattice equals the annihilator of the compact orbit
hull, and deleting any finite initial time interval leaves the same hull.
Every positive summable probe is continuous on that hull, so its recurrent
range is a compact interval; endpoints occur as limits, every interior value
is hit exactly, and compactness gives uniformly bounded waiting gaps.

For any two countable self-adjoint pure-point generators, one arbitrarily weak
self-adjoint trace-class separable direction can be chosen maximally
projective on every finite restriction.  Outside one countable coupling set,
the affine sum preserves the original maximal domain and self-adjointness,
has every finite temporal rank, has the full countable projective torus as
every future-tail hull, and has no nonzero exact global return.  Nevertheless
every positive summable probe retains complete syndetic recurrence.  The
finite-window all-probe difference remains bounded by
$2|t||g|\|W\|\|O\|$.

\section{Generic recurrence on a static qudit path}

Let $P_n$ be the open path and let $q\ge2$.  Its physical Hilbert space
has dimension $D=q^n$.  Let $\mathcal L_{P_n,q}^0$ be the trace-zero real
subspace spanned by arbitrary Hermitian onsite terms and finite sums of
Hermitian edge tensors.

\subsection{One local witness}

The proof constructs one simple maximally projective Hamiltonian inductively.
For two sites the edge algebra is the full Hermitian algebra.  A generic
eigenbasis makes the endpoint-observable transition graph complete.  Suppose
a chain has maximal projective energy and a connected non-bipartite endpoint
transition graph.  Attach one $q$-level site with separated diagonal
shifts.  Small edge perturbations activate a spanning family of cross-block
endpoint transitions.  For $q\ge3$, the local complete graph contains a
triangle.  For $q=2$, paired first-order cross- and same-layer transitions
create a triangle.  Simplicity, connectivity and non-bipartiteness persist
on an open set.

For an analytic eigenbranch relation
$\sum_{\alpha,r}c_{\alpha,r}\lambda_{\alpha,r}\equiv0$, vary one edge
matrix element and the shift difference $z=a_r-a_s$.  Second-order
perturbation gives
\[
 \sum_{\alpha,\beta}
 \frac{(c_{\alpha,r}-c_{\beta,s})
 |\langle\psi_\alpha,A\psi_\beta\rangle|^2}
 {z+E_\alpha-E_\beta}=0.
\]
Distinct old gaps make the poles distinct.  Their residues force equality
of coefficients along the categorical product of the endpoint graph with
$K_q$.  Connectedness plus an odd cycle makes this product connected, so
all coefficients are equal.  Only the trace relation survives.  Subtracting
the scalar trace part preserves projective maximality and lands in
$\mathcal L_{P_n,q}^0$.

\subsection{Genericity and recurrence consequences}

For every nonconstant balanced rational coefficient vector, multiply its
eigenvalue linear form over all label permutations.  The product is a
symmetric polynomial, hence a polynomial in characteristic coefficients and
therefore in local Hamiltonian parameters.  The witness proves it is not the
zero polynomial.  Its real zero set is null and nowhere dense.  The
discriminant is handled identically.  Removing the countable union over
rational relations yields a dense full-measure $G_\delta$.

\begin{theorem}[Exact path recurrence law]
For every $n,q\ge2$, the maximal set in
$\mathcal L_{P_n,q}^0$ is dense, residual and full Lebesgue measure.  Every
member has no nonzero balanced integer relation and anchored temporal rank
$q^n-1$.  For $0<\varepsilon<2$,
\[
 \lim_{T\to\infty}\frac1T\left|\left\{0\le t\le T:
 \max_{j\ne j_0}|\ee^{-\ii t(E_j-E_{j_0})}-1|<\varepsilon
 \right\}\right|
 =\left(\frac{2\arcsin(\varepsilon/2)}\pi\right)^{q^n-1}.
\]
No nonzero exact projective return exists.
\end{theorem}

\begin{proof}
The preceding algebraic avoidance proves maximality.  Its balanced-relation
definition gives zero kernel; rational rank of the anchored gaps is therefore
$D-1$.  Kronecker--Weyl equidistribution identifies time frequency with
Haar volume.  Each coordinate target arc has normalized length
$2\arcsin(\varepsilon/2)/\pi$, and independence multiplies the $D-1$
factors.  With at least three levels, two independent anchored frequencies
cannot both have a common nonzero exact period.
\end{proof}

\section{Reverse-water target information and deadlines}
\label{app:water}

Let $w_j>0$, $\sum_jw_j=1$, and $0<\eta\le1$.  For pivot $k$,
\[
 \Hwater_k=\min\left\{\frac12\sum_{j\ne k}\log\frac{w_j}{D_j}:
 0<D_j\le w_j,\ \sum_{j\ne k}D_j\le\eta\right\}.
\]
If $\eta\ge1-w_k$, its value is zero.  Otherwise a unique water level
$\lambda_k>0$ satisfies
$\sum_{j\ne k}\min\{w_j,\lambda_k^{-1}\}=\eta$, and
\[
 D_j^*=\min\{w_j,\lambda_k^{-1}\},\qquad
 \Hwater_k=\frac12\sum_{j\ne k}\log_+(\lambda_kw_j).
\]
Convexity and the Karush--Kuhn--Tucker conditions prove global optimality;
compactness after closing zero-cost inactive coordinates gives attainment.
Minimize over $k$ to obtain $\Hwater_{\rm proj}$.

\subsection{Weighted phase variance}

For phases $z_j\in\T$,
\[
 1-\left|\sum_jw_jz_j\right|^2
 =\frac12\sum_{i,j}w_iw_j|z_i-z_j|^2.
\]
This exact variance identity supplies both the product-box inclusion and the
observer interpretation.

\subsection{Target-volume sandwich}

Fix a pivot and require selected phase coordinates to lie in symmetric arcs
with squared chordal errors proportional to $D_j/w_j$.  The variance
identity shows this product box lies inside the fidelity target.  Its Haar
volume is bounded below by
\[
 \left(\frac{\sqrt2}{\pi}\right)^{r-1}
 \exp\{-\Hwater_{\rm proj}\}.
\]
Conversely, resolve the target by a pivot of maximal weighted alignment and
bound it by a quadratic ellipsoid/product envelope.  Summing over $r$
pivots gives
\[
 \mu_w(\eta)\le
 r\exp\left\{\frac\pi4(r-1)-\Hwater_{\rm proj}\right\}.
\]
Taking minus logarithms proves the two-sided target-information law and the
absolute and relative remainder bounds.  The constructive inner box also
proves strict target-volume positivity.

\subsection{Universal and selected-support clocks}

Round each coordinate resolution to an integer
$Q_j=\lceil\pi\sqrt{2w_j/D_j}\rceil$.  The anisotropic pigeonhole principle
applied to the first $\prod Q_j+1$ phase packets gives two in the same box;
their time difference is a positive return.  Minimizing and comparing the
rounded product with the convex objective proves
\[
 \Hwater_{\rm proj}-\frac\pi4(r-1)-\log(2r)
 <\log T_w
 \le\Hwater_{\rm proj}+(r-1)\log(2\pi\sqrt2).
\]
Applying the same argument to the shifted block of times
$M+1,\ldots,(M+1)T_w$ gives a return after every delay $M$.

For a selected set $S$ and discarded tail mass $\tau_S$, the inequality
\[
 2\tau_S+\sum_{j\in S}D_j\le\eta
\]
is sufficient.  Only $S$ enters the rounded product deadline; the tail pays
twice its weight.  This is an exact certificate for the selected set, not a
claim of joint discrete optimality.

For a phase vector $\omega$, multiplication by any integer $n$ preserves
Haar measure.  A union bound over the first $N$ ticks and the pivot-resolved
target envelope proves the finite-horizon Haar hit bound.  The complementary
mass gives an explicit $1-\alpha$ avoidance quantile whenever the budget is
at most $\alpha$.  Correlations between power-map targets preclude an
equality claim.

Independent onsite phases in a particle-number-preserving local XX spin
Hamiltonian give a diagonal single-excitation block.  The state
$\sum_j\sqrt{w_j}|j\rangle$ has return amplitude
$\sum_jw_jz_j$.  Thus the projective target and its Haar volume transfer
exactly to a local spin ensemble.  A binary construction proves that equal
hard effective support at one operational radius can have different water
entropy and hence different deadline bands.

\section{Observer coherence, operational probes and countable extension}

Let $B_a$ be a finite family of observables, let $b_a(i,j)$ denote the
matrix block from energy sector $j$ to sector $i$, and write
$|b_a(i,j)|$ for its Hilbert--Schmidt norm.  For weights
$\lambda_a\ge0$, put
\[
 W_{ij}=\sum_a\lambda_a|b_a(i,j)|^2,
 \qquad
 \mathcal E(z)=\sum_{i,j}W_{ij}|z_i-z_j|^2.
\]
The coherence graph has edge $ij$ exactly when $W_{ij}>0$.

\begin{theorem}[Finite observer capstone]
For Hamiltonian phases $z_i(t)=\ee^{-\ii tE_i}$, total squared observer
displacement equals $\mathcal E(z(t))$.  This quantity is zero if and only
if phases agree on every edge, equivalently on every connected component.
The range rank of the graph coboundary, and hence visible projective phase
capacity, is $r-c$, where $c$ is the number of components.
\end{theorem}

The identity follows by blockwise expansion.  The zero law uses
nonnegativity.  The rank formula follows because the kernel of the graph
coboundary consists of functions constant on components.

Synchronizing a spanning forest yields a recurrence by
$(Q+1)^{r-c}$ ticks with explicit squared-distance error.  Independent hard
frequencies on a rank-$r-c$ graph show that no smaller universal exponent is
possible.  Cauchy--Schwarz controls every Hilbert--Schmidt normalized linear
probe by total displacement.  An operator contraction has Hilbert--Schmidt
energy at most matrix dimension, giving the square-root dimension conversion;
the identity contraction attains the general factor.  A triangle inequality
adds two calibration errors.  If the resulting tick horizon is below the
certified coherent window, the recurrence is a device-feasible theorem.

For countable pure-point blocks, the same phase-energy identity holds exactly
when the nonnegative edge series is summable.  Truncating to a finite head,
synchronizing it and bounding the tail proves quantitative recurrence and the
zero law.  In the Hamiltonian specialization, zero means every positive
nonzero observer block is phase locked.

\subsection{Equal local information, opposite global recurrence}

There exist two explicit full-rank qutrit states with a common spectrum and
positive spectral weights such that their weighted variance, nonzero
Hilbert--Schmidt speed and nonzero symmetric-logarithmic-derivative quantum
Fisher information agree.  For the same tick protocol, the first returns at
every tick.  The second uses an irrational frequency, has an
inverse-resolution lower bound before return and a matching constructive
upper bound.  Therefore these standard local information quantities do not
determine global recurrence arithmetic.

\section{Complete finite-record ambiguity}

Fix $m\in\N$, $T>0$, $0\le H<T$, and $\delta>0$.  Choose an odd
$n>m$ so large that
\[
 2\left(\sin\frac{\pi H}{2T}\right)^n<\delta.
\]
Let $E_j=\pi j/T$ for $j=0,\ldots,n$, and define
\[
 q_j=2^{-n}\binom{n}{j},\qquad
 r_j=q_j(1+(-1)^{n-j}),\qquad
 w_{p,j}=(1-\sqrt p)q_j+\sqrt p\,r_j.
\]

\subsection{Positive mass and generating functions}

The binomial theorem gives $\sum_jq_j=1$.  The alternating binomial sum
gives $\sum_jr_j=1$, and all weights are nonnegative.  Hence every $w_p$
is a probability measure.  Its amplitude is
\[
 A_p(t)=\left(\frac{1+\ee^{\ii\pi t/T}}2\right)^n
 +\sqrt p\left(\frac{\ee^{\ii\pi t/T}-1}2\right)^n.
\]
At $t=T$, $A_p(T)=(-1)^n\sqrt p$, so $P_p(T)=p$.  Moreover
\[
 |A_p(t)-A_0(t)|=\sqrt p
 \left|\sin\frac{\pi t}{2T}\right|^n,
\]
and because both amplitudes have norm at most one,
\[
 |P_p(t)-P_0(t)|\le2|A_p(t)-A_0(t)|.
\]
This proves uniform continuous-window agreement simultaneously for every
$p\in[0,1]$.

\subsection{Shared Hamiltonian and uniform-field realization}

On $\C^{n+1}$, let $K=\operatorname{diag}(E_0,\ldots,E_n)$ and
$|\psi_p\rangle=\sum_j\sqrt{w_{p,j}}|j\rangle$.  One nonnegative
self-adjoint Hamiltonian therefore realizes all probabilities by varying
normalized preparations.

On $n$ qubits, the permutation-invariant uniform field
\[
 H_{\rm u}=\frac\pi T\sum_{a=1}^n|1\rangle\langle1|_a
\]
has symmetric Dicke eigenvectors of energies $E_j$.  The symmetric pure
state with Dicke-shell weights $w_{p,j}$ reproduces the same amplitude and
moments.  Thus one homogeneous noninteracting Hamiltonian gives the complete
closed-interval ambiguity when preparation may vary.

\subsection{Fixed preparation, projector and spectrum}

Fix $\phi=|0\rangle$.  Choose a unitary $U_p\phi=\psi_p$ and put
$K_p=U_p^*KU_p$.  All $K_p$ are nonnegative, self-adjoint and unitarily
equivalent.  Preparation $\phi$, return projector
$|\phi\rangle\langle\phi|$, spectrum and multiplicities are fixed, while
\[
 \langle\phi,\ee^{\ii tK_p}\phi\rangle=A_p(t).
\]

For $k<n$, the finite-difference identity
\[
 \sum_{j=0}^n(-1)^{n-j}{n\choose j}j^k=0
\]
implies that all $w_p$ have the same energy moments.  This proves the fixed
protocol, arbitrary moment and every-future clauses.

\subsection{Nearest-neighbour path and endpoint rigidity}

For $p<1$, all $w_{p,j}$ are positive.  The Krylov vectors
$\psi_p,K\psi_p,\ldots,K^n\psi_p$ are linearly independent by the
Vandermonde determinant.  Gram--Schmidt yields a basis starting at
$\psi_p$, and self-adjointness makes $K$ tridiagonal in that basis.
Strictly positive spectral weights make every forward Jacobi coupling
nonzero.

At $p=1$, only one parity of energy levels is occupied.  Tridiagonalize the
occupied Krylov block and place the unused levels in a diagonal spectator
block.  This gives a tridiagonal representative with one zero interface
coupling.  Conversely, the endpoint of a connected tridiagonal path is
cyclic: the eigenvalue equation propagates a zero endpoint coordinate through
the path.  An unused eigenvector orthogonal to the fixed endpoint therefore
forces at least one zero coupling.  This proves the exact connectedness
dichotomy and the closed-interval nearest-neighbour realization.

\subsection{Occupation and all-sector fermionic lifts}

Every Hermitian tridiagonal one-particle matrix is the single-excitation block
of a Hermitian particle-number-preserving nearest-neighbour occupation
Hamiltonian, with diagonal onsite terms and adjacent hopping terms.  For the
canonical second quantization,
\[
 d\Gamma(K_p)(u_{j_1}\wedge\cdots\wedge u_{j_k})
 =\left(\sum_{\ell=1}^kE_{j_\ell}\right)
  (u_{j_1}\wedge\cdots\wedge u_{j_k}).
\]
The Slater determinants form a basis.  Thus every $p$ has the same complete
subset-sum spectrum in every fermion-number sector, including multiplicity,
while the recurrence ambiguity remains in the one-excitation protocol.

\subsection{Sharp support and forecast error}

On $r$ common distinct nodes, equality of moments of orders
$0,\ldots,r-1$ gives a square Vandermonde system for the weight difference;
its determinant is nonzero.  Hence distinct measures matching the first
$n$ moments require at least $n+1$ nodes.  The binomial/parity pair uses
exactly $n+1$, so the support threshold is sharp.

Any forecaster receives one input compatible with terminal values zero and
one.  For its real output $y$,
$\max\{|y|,|y-1|\}\ge1/2$.  Therefore every forecasting rule has error at
least $1/2$ on one compatible clock.

\section{Global data boundary}

Finite moments are not the only incomplete local data.  There is no
universal classifier of persistent mean recurrence from complete polynomial
local data on positive-energy spectra with all moments.  Likewise, no
classifier based on the complete perturbative jet of every bounded continuous
spectral observable, or of every perturbative spectral-form-factor order,
works uniformly over compact positive-spectrum families.  The separating
families agree to all specified local or perturbative orders while having
opposite global recurrence phases.

Compact support is the exact positive moment boundary used here.  If two
probability measures are carried by one compact interval and have the same
Hamburger moment sequence, polynomial density in the continuous functions on
that interval makes their integrals agree on every continuous function;
Riesz uniqueness gives equality of measures.

\section{Guide to the results}

\begin{longtable}{@{}>{\raggedright\arraybackslash}p{0.36\textwidth}
  >{\raggedright\arraybackslash}p{0.56\textwidth}@{}}
\toprule
Result & Where established \\
\midrule
\endfirsthead
\toprule
Result & Where established \\
\midrule
\endhead
Projective completion; relation hull; $\tau+\rho=D-1$ & Main Law II;
Appendix Sections C and D \\
Singular wall, sinc profile, dense ranks, bounded-horizon indistinguishability
& Main Theorem 3; Appendix Section D \\
Protected affine relation and complete weak separable conversion & Main
Theorem 3; Appendix Sections D.1--D.2 \\
Countable relation hull and infinite separable recurrence ray & Main Law II;
Appendix Section D.3 \\
Replica height, finite-grid readout and Nyquist sharpness & Main Law II;
Appendix Section C.1 \\
Finite contraction threshold and coupled post-burn-in deadline & Main Law I;
Appendix Section B.1 \\
Stationary point-spectrum scaling and collision hierarchy & Main Law I;
Appendix Section B.2 \\
Countable pure-point recurrence, cyclic exact return and atomless boundary &
Main Law I; Appendix Section B.2 \\
Generic static nearest-neighbour qudit maximality and exact return frequency &
Main Theorem 4; Appendix Section E \\
Weighted projective variance and reverse-water global optimizer & Main Law
III; Appendix Sections F and F.1 \\
Product-box lower law, pivot upper law and two-sided target information & Main
Theorem 5; Appendix Section F.2 \\
Universal deadline, post-delay return, selected support and Haar hit quantile
& Main Theorem 5; Appendix Section F.3 \\
Exact onsite local-spin phase transfer and binary hard-support separation &
Main Law III; Appendix Section F.3 \\
Finite observer phase energy, zero set and graph capacity & Main Theorem 6;
Appendix Section G \\
Sharp graph exponent, probe, calibration and device-window bounds & Main
Theorem 6; Appendix Section G \\
Countable observer recurrence and zero law & Main Law III; Appendix Section
G \\
Equal qutrit local information with opposite global recurrence & Main Law
III; Appendix Section G.1 \\
Positive binomial/parity measures and amplitude identities & Main Law IV;
Appendix Section H.1 \\
One Hamiltonian, uniform-field qubit and fixed-protocol isospectral forms &
Main Theorem 7; Appendix Sections H.2--H.3 \\
Krylov path realization, perfect-return block and forced disconnection & Main
Theorem 7; Appendix Section H.4 \\
Occupation embedding and complete all-sector fermionic spectra & Main Theorem
7; Appendix Section H.5 \\
Common moments, exponential hiding and every exact future & Main Theorem 7;
Appendix Sections H.1--H.3 \\
Vandermonde support threshold and one-half forecast lower bound & Main Law IV;
Appendix Section H.6 \\
Complete local/perturbative nonidentifiability and compact moment determinacy &
Main Law IV; Appendix Section I \\
\bottomrule
\end{longtable}
